\begin{blocksection}
\question Write \lstinline{keep_if}, which takes a linked list of integers \lstinline{lnk} and a predicate \lstinline{pred}, and returns a new linked list of only the elements for which \lstinline{pred} returns true, in their original order.

\begin{lstlisting}
from collections.abc import Callable
def keep_if(
    lnk: LinkedList[int], pred: Callable[[int], bool]
) -> LinkedList[int]:
    """
    >>> lnk = Link(1, Link(2, Link(3, Link(4))))
    >>> print(keep_if(lnk, lambda x: x % 2 == 0))
    (2 4)
    >>> print(keep_if(lnk, lambda x: x > 0))
    (1 2 3 4)
    >>> keep_if(lnk, lambda x: x > 10)
    ()
    """
    if ______________________________________:

        return ______________________________

    rest = ____________________________________

    if ______________________________________:

        return ______________________________

    return ____________________________________
\end{lstlisting}

\begin{solution}[1in]
\begin{lstlisting}
def keep_if(
    lnk: LinkedList[int], pred: Callable[[int], bool]
) -> LinkedList[int]:
    if not isinstance(lnk, Link):
        return ()
    rest = keep_if(lnk.rest, pred)
    if pred(lnk.first):
        return Link(lnk.first, rest)
    return rest
\end{lstlisting}
\end{solution}

\end{blocksection}

\begin{questionmeta}
\textbf{Teaching Tips}\par\nopagebreak[4]
Suggested Time: 8 min; Difficulty: Easy--Medium
\nopagebreak[4]
\begin{itemize}
    \item \textbf{Leap of faith:} \lstinline{keep_if(lnk.rest, pred)} is the filtered version of the rest. Only the first element is left to decide.
    \item \textbf{Two outcomes:} Keep \lstinline{first} (put it on the front of the filtered rest) or drop it (return the filtered rest unchanged).
    \item \textbf{Connect to lists:} This is the linked-list version of a list comprehension with an \lstinline{if} clause.
\end{itemize}
\end{questionmeta}
